%! TEX root = master.tex
\chapter{1}{Cauchy's Integral Formula and Liouville's Theorem}

An analytic function is pinned down by its values on a closed curve: Cauchy's
integral formula recovers $f$ and all its derivatives inside from an integral
around the edge. Bounding that integral gives Cauchy's estimates, and from
them Liouville's theorem and the polynomial growth results.

\section{Cauchy's theorem}\label{sec:cauchy-theorem}

%! TEX root = ../master.tex
\begin{theorem}[$ML$-inequality]\label{thm:ml-inequality}
  If $f$ is continuous on a contour $\Gamma$ of length $L$ and
  $\abs{f(z)} \leq M$ for all $z$ on $\Gamma$, then
  \[
    \abs*{\int_\Gamma f(z) \dd z} \leq ML.
  \]
  \thmsource{\cite[ch.~4]{saff-snider}}
\end{theorem}

%! TEX root = ../master.tex
\begin{theorem}[Cauchy's integral theorem]\label{thm:cauchy-theorem}
  If $f$ is analytic in a simply connected domain $D$ and $\Gamma$ is any
  closed contour in $D$, then
  \[
    \oint_\Gamma f(z) \dd z = 0.
  \]
  \thmsource{\cite[ch.~4]{saff-snider}}
\end{theorem}

%! TEX root = ../master.tex
\begin{theorem}[Deformation of contours]\label{thm:deformation}
  Let $C$ be a simple closed positively oriented contour, and let
  $C_1, \dots, C_n$ be simple closed positively oriented contours inside $C$
  whose interiors are pairwise disjoint. If $f$ is analytic on a domain
  containing all these contours and the region inside $C$ but outside every
  $C_k$, then
  \[
    \oint_C f(z) \dd z = \sum_{k=1}^{n} \oint_{C_k} f(z) \dd z.
  \]
  \thmsource{\cite[ch.~4]{saff-snider}}
\end{theorem}


\section{Cauchy's integral formula}\label{sec:cif}

%! TEX root = ../master.tex
\begin{theorem}[Cauchy's integral formula]\label{thm:cif}
  Let $f$ be analytic in a simply connected domain $D$, and let $\Gamma$ be a
  simple closed positively oriented contour in $D$. For every $z_0$ inside
  $\Gamma$,
  \[
    f(z_0) = \frac{1}{2\pi i} \oint_\Gamma \frac{f(z)}{z - z_0} \dd z.
  \]
  \thmsource{\cite[ch.~4]{saff-snider}}
\end{theorem}

\begin{proof}
  Let $C_r$ be the circle $\abs{z - z_0} = r$, with $r$ small enough that
  $C_r$ lies inside $\Gamma$. The integrand is analytic everywhere except at
  $z_0$, so by \cref{thm:deformation}
  \[
    \oint_\Gamma \frac{f(z)}{z - z_0} \dd z
    = \oint_{C_r} \frac{f(z)}{z - z_0} \dd z
    = f(z_0) \oint_{C_r} \frac{\dd z}{z - z_0}
      + \oint_{C_r} \frac{f(z) - f(z_0)}{z - z_0} \dd z.
  \]
  Parametrising $z = z_0 + re^{it}$, $0 \leq t \leq 2\pi$, the first integral
  is $\int_0^{2\pi} i \dd t = 2\pi i$. Let $E_r$ be the second. On $C_r$ its
  integrand is at most $\max_{C_r} \abs{f(z) - f(z_0)} / r$ in modulus, so by
  \cref{thm:ml-inequality}
  \[
    \abs{E_r} \leq 2\pi \max_{\abs{z - z_0} = r} \abs{f(z) - f(z_0)}
    \longrightarrow 0
    \quad \text{as } r \to 0,
  \]
  by continuity of $f$ at $z_0$. But $E_r = \oint_\Gamma \frac{f(z)}{z - z_0}
  \dd z - 2\pi i f(z_0)$ does not depend on $r$, so $E_r = 0$.
\end{proof}

%! TEX root = ../master.tex
\begin{theorem}[Cauchy's integral formula for derivatives]\label{thm:cif-derivatives}
  Under the hypotheses of \cref{thm:cif}, $f$ has derivatives of every order
  in $D$, each of them analytic in $D$, and for every $z_0$ inside $\Gamma$
  \[
    f^{(n)}(z_0) = \frac{n!}{2\pi i} \oint_\Gamma \frac{f(z)}{(z - z_0)^{n+1}} \dd z
    \qquad (n = 0, 1, 2, \dots).
  \]
  \thmsource{\cite[ch.~4]{saff-snider}}
\end{theorem}


\begin{exercise}\label{ex:cos-over-square}
  Let $f(z) = \dfrac{\cos z}{(z^2 + 1)^2}$.
  \begin{enumerate}[label=(\alph*)]
    \item Find the singularities of $f$.
    \item Show that the integral of $f$ around $\abs{z} = 3$ is the sum of
      its integrals around small circles centred at the singularities.
    \item Evaluate $\oint_{\abs{z} = 3} f(z) \dd z$.
  \end{enumerate}
\end{exercise}

\begin{solution}
  \emph{Singularities.} $\cos z = \frac12(e^{iz} + e^{-iz})$ is entire and
  $(z^2 + 1)^2$ is a polynomial, so $f$ is analytic except where
  $(z^2 + 1)^2 = 0$, that is, at $z = \pm i$.

  \emph{Decomposition.} Let $C_1$ and $C_2$ be the positively oriented
  circles $\abs{z - i} = \frac12$ and $\abs{z + i} = \frac12$. They lie inside
  $\abs{z} = 3$, their interiors are disjoint, and $f$ is analytic everywhere
  else. By \cref{thm:deformation},
  \[
    \oint_{\abs{z} = 3} f(z) \dd z
    = \oint_{C_1} f(z) \dd z + \oint_{C_2} f(z) \dd z.
  \]

  \emph{Around $i$.} Write $f(z) = g_1(z) / (z - i)^2$ with
  $g_1(z) = \cos z / (z + i)^2$, which is analytic on and inside $C_1$. By
  \cref{thm:cif-derivatives} with $n = 1$, the integral is $2\pi i \, g_1'(i)$.
  Now
  \[
    g_1'(z) = \frac{-(z + i) \sin z - 2\cos z}{(z + i)^3},
  \]
  and $\sin i = i \sinh 1$, $\cos i = \cosh 1$, so
  \[
    g_1'(i)
    = \frac{-2i \cdot i \sinh 1 - 2\cosh 1}{(2i)^3}
    = \frac{2(\sinh 1 - \cosh 1)}{-8i}
    = \frac{-2e^{-1}}{-8i}
    = \frac{1}{4ie}.
  \]
  Hence $\oint_{C_1} f(z) \dd z = 2\pi i \cdot \frac{1}{4ie} = \frac{\pi}{2e}$.

  \emph{Around $-i$.} Likewise, with $g_2(z) = \cos z / (z - i)^2$,
  \[
    g_2'(-i)
    = \frac{-(-2i)(-i \sinh 1) - 2\cosh 1}{(-2i)^3}
    = \frac{2(\sinh 1 - \cosh 1)}{8i}
    = -\frac{1}{4ie},
  \]
  so $\oint_{C_2} f(z) \dd z = -\frac{\pi}{2e}$.

  \emph{Total.}
  \[
    \oint_{\abs{z} = 3} \frac{\cos z}{(z^2 + 1)^2} \dd z
    = \frac{\pi}{2e} - \frac{\pi}{2e} = 0. \qedhere
  \]
\end{solution}

\begin{remark}
  The answer $0$ can be predicted without computing. $f$ is even, and the
  substitution $z = -w$ maps the circle $\abs{z} = 3$ onto itself with the
  same orientation, so
  $\oint f(z) \dd z = \oint f(-w)(-\dd w) = -\oint f(w) \dd w$.
\end{remark}

\begin{exercise}\label{ex:two-integrals}
  Let $\Gamma$ be a simple closed positively oriented contour, and let $f$ be
  analytic inside and on $\Gamma$. For $z_0$ not on $\Gamma$:
  \begin{enumerate}[label=(\alph*)]
    \item evaluate $\oint_\Gamma \frac{f'(z)}{z - z_0} \dd z$ by Cauchy's
      integral formula;
    \item evaluate $\oint_\Gamma \frac{f(z)}{(z - z_0)^2} \dd z$ by the
      formula for derivatives;
    \item compare the two;
    \item say what goes wrong when $z_0$ lies on $\Gamma$.
  \end{enumerate}
\end{exercise}

\begin{solution}
  By \cref{thm:cif-derivatives}, $f'$ is analytic wherever $f$ is.

  \emph{$z_0$ inside $\Gamma$.} \Cref{thm:cif} applied to $f'$, and
  \cref{thm:cif-derivatives} with $n = 1$ applied to $f$, give
  \[
    \oint_\Gamma \frac{f'(z)}{z - z_0} \dd z = 2\pi i f'(z_0)
    = \oint_\Gamma \frac{f(z)}{(z - z_0)^2} \dd z.
  \]

  \emph{$z_0$ outside $\Gamma$.} Both integrands are analytic on a simply
  connected domain containing $\Gamma$ and its inside but not $z_0$, so both
  integrals are $0$ by \cref{thm:cauchy-theorem}.

  \emph{Comparison.} In either case the two integrals are equal.

  \emph{$z_0$ on $\Gamma$.} Both integrands are unbounded near $z_0$, so they
  are not continuous on $\Gamma$. Neither formula applies, and the integrals
  are not even defined as ordinary contour integrals; they would have to be
  treated as improper integrals.
\end{solution}

\begin{remark}
  The equality has a one-line explanation. Off $z_0$,
  \[
    \frac{\dd}{\dd z} \frac{f(z)}{z - z_0}
    = \frac{f'(z)}{z - z_0} - \frac{f(z)}{(z - z_0)^2},
  \]
  and the integral of a derivative around a closed contour is $0$.
\end{remark}

\section{Cauchy's estimates and Liouville's theorem}\label{sec:liouville}

%! TEX root = ../master.tex
\begin{theorem}[Cauchy's estimates]\label{thm:cauchy-estimates}
  Let $f$ be analytic inside and on the circle $C_R \colon \abs{z - z_0} = R$,
  with $\abs{f(z)} \leq M$ on $C_R$. Then
  \[
    \abs{f^{(n)}(z_0)} \leq \frac{n! \, M}{R^n}
    \qquad (n = 0, 1, 2, \dots).
  \]
  \thmsource{\cite[ch.~4]{saff-snider}}
\end{theorem}

\begin{proof}
  On $C_R$ the integrand in \cref{thm:cif-derivatives} satisfies
  $\abs{f(z) / (z - z_0)^{n+1}} \leq M / R^{n+1}$, and $C_R$ has length
  $2\pi R$. By \cref{thm:ml-inequality},
  \[
    \abs{f^{(n)}(z_0)}
    \leq \frac{n!}{2\pi} \cdot \frac{M}{R^{n+1}} \cdot 2\pi R
    = \frac{n! \, M}{R^n}. \qedhere
  \]
\end{proof}

%! TEX root = ../master.tex
\begin{theorem}[Liouville's theorem]\label{thm:liouville}
  A bounded entire function is constant.
  \thmsource{\cite[ch.~4]{saff-snider}}
\end{theorem}

\begin{proof}
  Suppose $\abs{f} \leq M$ on $\C$. Fix $z_0$. For every $R > 0$,
  \cref{thm:cauchy-estimates} with $n = 1$ gives
  $\abs{f'(z_0)} \leq M / R$. Letting $R \to \infty$ gives $f'(z_0) = 0$.
  So $f' \equiv 0$ on the connected set $\C$, and $f$ is constant.
\end{proof}


\begin{exercise}\label{ex:bounded-imaginary-part}
  Let $f$ be entire with $A \leq \operatorname{Im} f(z) \leq B$ for all $z$.
  Show that $e^{if}$ is bounded and entire, and deduce that $f$ is constant.
\end{exercise}

\begin{proof}
  Write $f = u + iv$, so $if = -v + iu$ and
  \[
    \abs{e^{if}} = \abs{e^{-v} e^{iu}} = e^{-v}.
  \]
  Since $A \leq v \leq B$, this gives $e^{-B} \leq \abs{e^{if}} \leq e^{-A}$,
  so $e^{if}$ is bounded. It is entire because it is a composition of the
  entire functions $if$ and $\exp$, with
  $\frac{\dd}{\dd z} e^{if} = if' e^{if}$ by the chain rule.

  By \cref{thm:liouville}, $e^{if} \equiv c$ for some constant $c$, and
  $c \neq 0$ because the exponential never vanishes. Differentiating,
  $if'(z) e^{if(z)} = 0$, so $f' \equiv 0$ and $f$ is constant.
\end{proof}

\begin{exercise}\label{ex:bounded-seventh-derivative}
  Let $f$ be entire with $f^{(7)}$ bounded on $\C$. What kind of function is
  $f$? Generalise, and decide whether the same holds for real functions.
\end{exercise}

\begin{solution}
  By \cref{thm:cif-derivatives}, $f^{(7)}$ is entire, and it is bounded, so it
  is a constant $c$ by \cref{thm:liouville}. Hence $f^{(n)} \equiv 0$ for
  $n \geq 8$. By Taylor's theorem (\cref{thm:taylor}) the series of $f$ about
  $0$ converges to $f$ on all of $\C$, and its coefficients
  $f^{(n)}(0) / n!$ vanish for $n \geq 8$. So
  \[
    f(z) = \sum_{n=0}^{7} \frac{f^{(n)}(0)}{n!} z^n,
  \]
  a polynomial of degree at most $7$, and exactly $7$ when $c \neq 0$.

  \emph{Generalisation.} The same argument shows that if $f$ is entire and
  $f^{(k)}$ is bounded, then $f$ is a polynomial of degree at most $k$.

  \emph{Real functions.} This fails on $\R$. For $g(x) = \sin x$, every
  derivative is $\pm \sin x$ or $\pm \cos x$, so $g^{(7)}(x) = -\cos x$ is
  bounded, yet $\sin x$ is not a polynomial.
\end{solution}

\begin{remark}
  Applying \cref{thm:cauchy-estimates} directly to $g = f^{(7)}$ gives
  $\abs{f^{(8)}(z_0)} \leq M / R$ for every $R$, which is the same argument
  that proves Liouville's theorem.
\end{remark}

\section{The maximum modulus principle}\label{sec:maximum-modulus}

%! TEX root = ../master.tex
\begin{theorem}[Maximum modulus principle]\label{thm:maximum-modulus}
  Let $D$ be a bounded domain, and let $f$ be continuous on the closure of $D$
  and analytic in $D$. Then $\abs{f}$ attains its maximum over the closure on
  the boundary of $D$. If $\abs{f}$ also attains it at a point of $D$, then $f$
  is constant.
  \thmsource{\cite[ch.~4]{saff-snider}}
\end{theorem}


\begin{exercise}\label{ex:max-on-disc}
  Let $f(z) = \left(z - \frac12\right)\left(z - \frac14\right)$ on the closed
  unit disc. Explain why $\max_{\abs{z} \leq 1} \abs{f(z)}$ is attained on
  $\abs{z} = 1$, compute $\abs{f}$ there, and find the maximum.
\end{exercise}

\begin{solution}
  \emph{Where.} $f$ is a polynomial, so it is analytic on the open unit disc
  and continuous on its closure. By \cref{thm:maximum-modulus}, the maximum of
  $\abs{f}$ is attained on $\abs{z} = 1$.

  \emph{On the circle.} Write $z = e^{i\theta}$ and $a = \cos\theta$. For real
  $c$, $\abs{z - c}^2 = (z - c)(\bar z - c) = 1 - 2ca + c^2$, so
  \[
    \abs{f(z)}^2
    = \left(\tfrac54 - a\right)\left(\tfrac{17}{16} - \tfrac12 a\right)
    = \tfrac12 a^2 - \tfrac{27}{16} a + \tfrac{85}{64}.
  \]

  \emph{Maximum.} The derivative in $a$ is $a - \frac{27}{16} < 0$ on
  $[-1, 1]$, so $\abs{f}^2$ is largest at $a = -1$, that is, at $z = -1$, where
  it equals $\frac{225}{64}$. Hence
  \[
    \max_{\abs{z} \leq 1} \abs{f(z)} = \frac{15}{8}. \qedhere
  \]
\end{solution}

\begin{remark}
  The triangle inequality gets there faster: for $\abs{z} \leq 1$,
  $\abs{f(z)} \leq \left(1 + \frac12\right)\left(1 + \frac14\right) = \frac{15}{8}$,
  with equality at $z = -1$.
\end{remark}
